\documentclass[12pt, twoside]{article}
%\documentclass[12pt, twoside]{article}
\usepackage[letterpaper, margin=1in, headsep=0.2in]{geometry}
\setlength{\headheight}{0.6in}
%\usepackage[english]{babel}
\usepackage[utf8]{inputenc}
\usepackage{microtype}
\usepackage{amsmath}
\usepackage{amssymb}
%\usepackage{amsfonts}
\usepackage[nomessages]{fp} %\FPeval{\var-name}{2*sin(pi/6)}
\usepackage{siunitx} %units in math. eg 20\milli\meter
\usepackage{yhmath} % for arcs, overparenth command
\usepackage{tikz} %graphics
\usetikzlibrary{quotes, angles, arrows, arrows.meta}
\usepackage{pgfplots}
\usepgfplotslibrary{fillbetween}
\pgfplotsset{compat=1.16}
\usepackage{graphicx} %consider setting \graphicspath{{images/}}
\usepackage{parskip} %no paragraph indent
\usepackage{enumitem}
\usepackage{multicol}
\usepackage{venndiagram}

\usepackage{fancyhdr}
\pagestyle{fancy}
\fancyhf{}
\renewcommand{\headrulewidth}{0pt} % disable the underline of the header
\raggedbottom
\hfuzz=2mm %suppresses overfull box warnings

\usepackage{hyperref}

\fancyhead[LE]{\thepage}
\fancyhead[RO]{Solutions}
\fancyhead[LO]{La Scuola d'Italia / Dr.\ Huson / IB Math 11 \\* 7 October 2026}

\begin{document}
\subsubsection*{1.10 Early Finishers: Sequences and Compound Interest --- Solutions}

\emph{Extension set, not scored: worked solutions only, no marks.
Calculator allowed. Answers are exact or correct to three significant
figures.}

\begin{enumerate}[itemsep=0.3cm]

%--- Problem 1: Arithmetic sequence from two conditions; count of terms below 500; arithmetic terms forming a geometric sequence ---
\item[1.]

$u_{3} + u_{8} = 48$, \ $u_{12} = 5\,u_{2}$.

\medskip

\textbf{(a)} Write each term as $u_{1} + (n - 1)d$:

$(u_{1} + 2d) + (u_{1} + 7d) = 48 \;\Rightarrow\; 2u_{1} + 9d = 48$

$u_{1} + 11d = 5(u_{1} + d) \;\Rightarrow\; 6d = 4u_{1}
\;\Rightarrow\; u_{1} = 1.5d$

Substitute: $2(1.5d) + 9d = 48 \;\Rightarrow\; 12d = 48
\;\Rightarrow\; d = 4$

$u_{1} = 1.5(4) = 6$

\emph{Or rearrange both conditions to $2u_{1} + 9d = 48$ and
$4u_{1} - 6d = 0$ and solve the system on the GDC. The sequence is
$6, 10, 14, \ldots$: $u_{3} + u_{8} = 14 + 34 = 48$ and
$u_{12} = 50 = 5 \times 10$.}

\medskip

\textbf{(b)} $6 + 4(n - 1) < 500 \;\Rightarrow\; 4(n - 1) < 494
\;\Rightarrow\; n - 1 < 123.5 \;\Rightarrow\; n < 124.5$

$u_{124} = 6 + 4(124 - 1) = 498$ and $u_{125} = 6 + 4(125 - 1) = 502$,
so $124$ terms are less than $500$.

\emph{The terms below $500$ are $u_{1}$ to $u_{124}$, so the count is the
largest whole $n$, $124$.}

\medskip

\textbf{(c)} The $1^{\text{st}}$, $2^{\text{nd}}$ and $4^{\text{th}}$
terms are $3$, \ $3 + d$, \ $3 + 3d$.

Geometric, so the ratios are equal:
$\dfrac{3 + d}{3} = \dfrac{3 + 3d}{3 + d}$

$(3 + d)^{2} = 3(3 + 3d) \;\Rightarrow\; 9 + 6d + d^{2} = 9 + 9d
\;\Rightarrow\; d^{2} - 3d = 0 \;\Rightarrow\; d(d - 3) = 0$

$d \neq 0$, so $d = 3$.

The terms are $3,\; 6,\; 12$, so $r = \dfrac{6}{3} = 2$.

\emph{$d = 0$ also solves the equation (the constant sequence
$3, 3, 3$), which is why the question excludes it.}

\newpage

%--- Problem 2: Compound interest -- doubling time; unknown rate; first year one investment overtakes another ---
\item[2.]

Sofia: $PV = 5000$, $4.8\%$ compounded monthly. \ Luca: $PV = 6000$,
$r\%$ compounded quarterly, $FV = 9000$ after $10$ years.

\medskip

\textbf{(a)} $FV = PV\left(1 + \dfrac{r}{100k}\right)^{kn}$, with
$k = 12$:

$5000\left(1 + \dfrac{4.8}{100(12)}\right)^{12n} > 10\,000$

\begin{tabular}{cc}
$n$ (years) & Value \\
\hline
$14$ & \$$9777.63$ \\
$15$ & \$$10\,257.42$ \\
\end{tabular}

$15$ years

\emph{Or TVM solver: $I\% = 4.8$, $PV = -5000$, $PMT = 0$, $FV = 10\,000$,
$P/Y = C/Y = 12$ gives $N = 173.6$ months $= 14.47$ years; the value is
first over \$$10\,000$ at a whole year after $15$ years. Show the adjacent
year: $14$ years is what shows $15$ is the \emph{least}.}

\medskip

\textbf{(b)} $6000\left(1 + \dfrac{r}{100(4)}\right)^{4(10)} = 9000$

$\left(1 + \dfrac{r}{400}\right)^{40} = 1.5$

$r = 400\left(1.5^{1/40} - 1\right) = 4.0752\ldots \approx 4.08$

\emph{TVM solver: $N = 40$, $PV = -6000$, $PMT = 0$, $FV = 9000$,
$P/Y = C/Y = 4$, solve for $I\% = 4.0752\ldots$}

\medskip

\textbf{(c)} Sofia: $S_{n} = 5000\left(1 + \dfrac{4.8}{100(12)}\right)^{12n}$
\qquad Luca: $L_{n} = 6000\left(1 + \dfrac{4.0752\ldots}{100(4)}\right)^{4n}$

\begin{tabular}{ccc}
$n$ (years) & Sofia & Luca \\
\hline
$24$ & \$$15\,786.26$ & \$$15\,877.07$ \\
$25$ & \$$16\,560.90$ & \$$16\,534.06$ \\
\end{tabular}

Year $25$

\emph{Use a GDC table of both expressions, with the full-precision $r$
stored from (b). Sofia starts with less but earns the higher effective
rate ($1.004^{12} = 1.0491$, $4.91\%$ a year, against Luca's
$1.5^{1/10} = 1.0414$, $4.14\%$ a year), so she catches up slowly. With
$r = 4.08$ rounded, Luca's values are \$$15\,894.92$ and \$$16\,553.42$:
the answer is still year $25$, but carry full precision to the end.}

\end{enumerate}
\end{document}
